\section{The problem}

One of the first problems we encounter with the use of the
\texttt{buffer.py} module\footnote{And obviously, any other parent
  version of \texttt{buffer.py}.} is that, if we don't use headphones,
the sound that comes out of our PC's (loud)speaker some time
(miliseconds) later reaches our mic(rophone), and more some time
later, that sound reaches our interlocutor (the ``far-end'' ... in the
system we are the ``near-end'') in the form of an echo (signal) of its
own voice, which is reproduced by his/her speaker, which can be
captured again (some time later) by his/her mic and sent it back to us
... and so on, generating a rather unpleasant feedback ``noise''.

To formalize this problem, let's define:
\begin{enumerate}
\item $s$ as the (analog) signal played by our
  speaker, % and that reaches our mic
\item $n$ the signal emited by
  the near-end person (that's me)\footnote{I.e.,
  the same signal that would be captured by our mic if I were using a
  headset.}, and % that reaches my mic
\item $m$ the (mixed) signal recorded by my microphone\footnote{Notice
that $m(t)$ is the analog audio signal that makes the membrane of our
microphone oscillate, and therefore, the signal that our sound card
records.}  ,
\end{enumerate}
where
\begin{equation}
  m(t) = \tilde{n}(t) + \tilde{s}(t),
  \label{eq:echo_problem}
\end{equation}
where $\tilde{\cdot}$ represents an approximation of $\cdot$. We have
here approximations because the signals are modified when they travel
through the air.

Our problem here is to minimize the
\href{https://en.wikipedia.org/wiki/Energy_(signal_processing)}{energy}
of $\tilde{s}$, i.e., to make
\begin{equation}
  \tilde{s}(t) = 0.
\end{equation}

\begin{comment}
\section{Echo cancellation}
In Intercom, the echo (signal) is a delayed copy of our voice,
$\mathbf{n}^{(\text{echo~delay})}$. With this idea in mind, ideally,
the echo could be removed using
\begin{equation}
  \tilde{\mathbf{n}}^{(i)} = \mathbf{m}^{(i)} - a\mathbf{n}^{(i-j)},
\end{equation}
where $i$ represents the $i$-th chunk, and $j$ the number of chunks
that are stored in the dejittering buffer.

Unfortunately, $\mathbf{n}$ is unknown, and the best approximation we
have for the echo at the $i$-th iteration, is
$\mathbf{m}^{(i-j)}$. For this reason, lets finally define the echo
cancellation algorithm (the signal that we should send after
cancelling the echo) as
\begin{equation}
  \tilde{\mathbf{n}}^{(i)} = \mathbf{m}^{(i)} - a\mathbf{m}^{(i-j)},
\end{equation}

\section{Modeling the echo signal}
Similar to how feedback suppression can be improved by using a filter
that models the listening environment, echo cancellation can be
improved by using a filter. Therefore, the signal that is finally sent
should be calculated as
\begin{equation}
  \tilde{n}^{(i)} = \mathbf{m}^{(i)} - (\mathbf{h}*\mathbf{m}^{(i-j)},
\end{equation}
where $\mathbf{h}$ represents the filter that models the listening
environment.

\section{Filter update using impulses}
If we consider that the listening environment can change, $\mathbf{h}$ should
adapt to these changes. At this point, again, we can use a sequence of
impulses to estimate $\mathbf{h}$, but considering that the effects of
the impulses will only be perceived after $j$ chunk-times. In this
case, we have that
\begin{equation}
  \tilde{n}^{(i)} = \mathbf{m}^{(i)} - (\mathbf{h}^{(i)}*\mathbf{m}^{(i-j)}.
\end{equation}

\end{comment}
